\chuong{TỔNG QUAN VỀ ĐƠN VỊ KIẾN TẬP}

Chương này giới thiệu đơn vị kiến tập theo ba lớp, đi từ rộng đến hẹp. Lớp ngoài cùng là công ty và lĩnh vực hoạt động, kế đến là bộ phận nơi nhóm làm việc, sau cùng là vị trí của nhóm trong bộ phận ấy. Qua ba lớp này, chương muốn cho thấy bài toán kiểm soát truy cập của AI Agent, vốn là chủ đề của các chương sau, xuất phát từ nhu cầu có thật của một công ty tài chính.

\subsection{Giới thiệu Công ty Tài chính Tổng hợp cổ phần Tín Việt}

\subsubsection{Lịch sử hình thành}

Đơn vị kiến tập là Công ty Tài chính Tổng hợp cổ phần Tín Việt (VietCredit General Finance Joint Stock Company, gọi tắt là VietCredit). Công ty được Ngân hàng Nhà nước Việt Nam (NHNN) cấp phép năm 2008 với tên Công ty Tài chính cổ phần Xi măng \parencite{vietcredit_bctn2025}. Mười năm sau, ngày 18/6/2018, NHNN cấp đổi giấy phép và công ty mang tên Công ty Tài chính cổ phần Tín Việt; thời điểm đó vốn điều lệ vào khoảng 605 tỷ đồng, và Tổng công ty Công nghiệp Xi măng Việt Nam (Vicem) là cổ đông lớn duy nhất \parencite{nhadautu_doiten2018}. Cũng từ năm 2018, công ty rẽ sang mảng tài chính tiêu dùng \parencite{vietcredit_bctn2025}. Chữ ``Tổng hợp'' trong tên hiện tại chỉ xuất hiện từ ngày 26/05/2025 \parencite{baodautu_doiten2025}.

Tính đến ngày 31/12/2025, VietCredit có vốn điều lệ khoảng 912 tỷ đồng và 454 cán bộ nhân viên. Trụ sở chính đặt tại Hà Nội, kèm một chi nhánh tại Thành phố Hồ Chí Minh, và cổ phiếu của công ty giao dịch trên UPCoM với mã TIN \parencite{vietcredit_bctn2025}.

\subsubsection{Lĩnh vực hoạt động}

Với tư cách công ty tài chính tổng hợp, VietCredit được huy động vốn, cấp tín dụng dưới các hình thức cho vay tiêu dùng, phát hành thẻ tín dụng và cho thuê tài chính, đồng thời cung cấp một số dịch vụ tài chính khác trong phạm vi giấy phép của NHNN \parencite{vietcredit_bctn2025}. Trên thực tế, trọng tâm là cho vay qua kênh số. Sản phẩm Tin Vay được nhúng vào ứng dụng của các đối tác như MoMo, Viettel Money, Zalo và Grab, nhờ vậy khách hàng vay hoàn toàn trực tuyến; song song có Tin Vay Biz dành cho hộ kinh doanh và nền tảng TIN GO cho vay mua xe điện. Đến cuối năm 2025, gần 80\% dư nợ đến từ cho vay qua công nghệ số \parencite{vietcredit_kqkd2025}. Quy mô vẫn tăng nhanh. Ngày 30/9/2026, dư nợ cho vay khách hàng chạm gần 20.928 tỷ đồng, cao hơn cùng kỳ trên 124\%, còn lợi nhuận trước thuế chín tháng đầu năm 2026 đạt 1.558,6 tỷ đồng \parencite{vnba_9t2026}.

Khi gần như toàn bộ hoạt động cho vay vận hành trên hệ thống, an toàn thông tin trở thành điều kiện để kinh doanh chứ không còn là việc riêng của bộ phận kỹ thuật. Báo cáo thường niên ghi rằng các hệ thống thông tin của công ty được phân loại từ cấp độ 1 đến cấp độ 3 theo quy định của NHNN, mỗi cấp độ đi kèm biện pháp bảo mật tương ứng; cùng lúc, công ty đặt chuyển đổi số, trong đó có trí tuệ nhân tạo và tự động hóa, làm trọng tâm chiến lược năm 2026 \parencite{vietcredit_bctn2025}. Hai định hướng ấy gặp nhau ở câu hỏi mà đề tài theo đuổi. Một AI Agent được phép thao tác trên những hệ thống đã phân cấp bảo mật thì phải chịu cùng khuôn khổ kiểm soát truy cập như một nhân viên.

\subsubsection{Quy trình làm việc}

Trong phạm vi đề tài, quy trình làm việc của nhóm có thể truy lại qua lịch sử kho mã. Mỗi thành phần được phát triển trên một nhánh riêng như \texttt{feat/keycloak}, \texttt{feat/rbac} hay \texttt{feat/grafana-allow-deny}, sau đó gộp vào nhánh \texttt{dev} để chạy thử chung, rồi mới đưa lên \texttt{main} bằng pull request. Từ 08/10 đến 04/11/2026, bảy pull request đã đi qua con đường này. Nhóm chỉ coi một thay đổi là hoàn tất khi kiểm thử tĩnh và kiểm thử đầu cuối liên quan cùng đạt, và đầu ra kiểm thử, sau khi loại bỏ secret cùng token, được lưu vào thư mục \path{evidence/}. Báo cáo cũng có nhánh riêng, được chia theo chương và hai thành viên đọc chéo cho nhau (Bảng~\ref{tab:phancong}). Công việc của từng tuần được ghi ở phần Lịch làm việc đầu báo cáo.

\subsection{Vị trí kiến tập và định hướng nghề nghiệp}

Đề tài đặt nhóm vào vùng giao nhau giữa kỹ thuật nền tảng (platform engineering), DevOps và observability, nhưng áp dụng cho một kiểu người dùng mới là AI Agent. Phần lớn thời gian không dành cho việc huấn luyện mô hình. Nhóm dồn sức vào dựng môi trường để agent chạy an toàn, cụ thể là viết manifest Kubernetes và script chạy lại được nhiều lần, cấu hình danh tính máy với máy trên Keycloak, soạn quy tắc phân quyền theo công cụ trên gateway và dựng dashboard kiểm toán trên Grafana. Bảng~\ref{tab:phancong} trình bày cách hai thành viên chia phần kỹ thuật.

Lý do chọn hướng này khá đơn giản. Kỹ năng nền tảng ít lệ thuộc vào một mô hình cụ thể, bởi mô hình ngôn ngữ có thể bị thay thế sau vài tháng trong khi yêu cầu cấp danh tính, phân quyền tối thiểu và ghi nhật ký kiểm toán vẫn giữ nguyên. Đơn vị cũng cần đúng những kỹ năng đó, vì tin tuyển dụng nhắc ở trên đòi hỏi Docker, Kubernetes, CI/CD và giám sát bằng Grafana \parencite{vietcredit_tuyendung_ds}. Qua đề tài, nhóm đã làm việc trực tiếp với tất cả các công cụ này, riêng CI/CD thì chưa. CI/CD, cùng với xoay vòng secret cho agent và cảnh báo chủ động dựa trên log, là những nội dung nhóm định học tiếp.
